\documentclass[12pt]{article}
\usepackage{amsmath, amssymb, amsthm}
\usepackage{geometry}
\geometry{margin=1in}

\title{CSE 400: Fundamentals of Probability in Computing\\
Tutorial 2 -- Lecture 14 \\ Exam-Oriented Lecture Scribe}
\date{}

\begin{document}
\maketitle

\section*{Tutorial Question 1}

\textbf{Problem.}
A point is chosen uniformly at random on a line segment of length $L$.  
Find the probability that the ratio of the shorter to the longer segment is less than $1/4$.

\textbf{Solution.}

Let $X$ denote the distance of the chosen point from the left endpoint.  
Since the point is chosen uniformly,
\[
X \sim \text{Uniform}(0,L).
\]

The two segment lengths are
\[
X \quad \text{and} \quad L-X.
\]

The ratio condition becomes
\[
\frac{\min(X,L-X)}{\max(X,L-X)} < \frac14.
\]

Consider symmetry about $L/2$.

\textbf{Case 1:} $X \le L/2$  
Shorter $=X$, longer $=L-X$,
\[
\frac{X}{L-X}<\frac14
\Rightarrow 4X<L-X
\Rightarrow 5X<L
\Rightarrow X<\frac{L}{5}.
\]

\textbf{Case 2:} $X>L/2$  
Shorter $=L-X$,
\[
\frac{L-X}{X}<\frac14
\Rightarrow 4(L-X)<X
\Rightarrow 4L<5X
\Rightarrow X>\frac{4L}{5}.
\]

Hence favorable region:
\[
X\in\left(0,\frac{L}{5}\right)\cup\left(\frac{4L}{5},L\right).
\]

Probability:
\[
P=\frac{L/5+L/5}{L}=\frac{2}{5}.
\]

\textbf{Final Result:} $\boxed{2/5}$.

\section*{Tutorial Question 2}

\textbf{Problem.}
White readings: $N(4,4)$, Black readings: $N(6,9)$.  
Observed value = 5. Black region fraction = $\alpha$.  
Find $\alpha$ such that error probabilities are equal.

\textbf{Solution.}

Equal error occurs when posterior probabilities are equal:
\[
P(\text{White}\mid 5)=P(\text{Black}\mid 5).
\]

Using Bayes rule proportionality,
\[
(1-\alpha)f_W(5)=\alpha f_B(5).
\]

Normal densities:
\[
f_W(5)=\frac{1}{\sqrt{8\pi}}
e^{-(5-4)^2/8},
\quad
f_B(5)=\frac{1}{\sqrt{18\pi}}
e^{-(5-6)^2/18}.
\]

Equating:
\[
(1-\alpha)\frac{1}{\sqrt{8\pi}}e^{-1/8}
=\alpha\frac{1}{\sqrt{18\pi}}e^{-1/18}.
\]

Solve for $\alpha$:
\[
\frac{1-\alpha}{\alpha}
=\sqrt{\frac{8}{18}}
e^{\,(-1/18+1/8)}.
\]

\textbf{Final Result:} $\alpha$ obtained from above equality.

\section*{Tutorial Question 3}

\subsection*{(a) Uniform Fires on $(0,A)$}

Minimize
\[
E|X-a|=\frac1A\int_0^A |x-a|dx.
\]

Split integral at $a$:
\[
\int_0^a (a-x)dx+\int_a^A(x-a)dx.
\]

Differentiate with respect to $a$ and set to zero:
\[
\frac{d}{da}E|X-a|=0
\Rightarrow a=\frac{A}{2}.
\]

\textbf{Result:} Station at midpoint.

\subsection*{(b) Exponential Fires}

$X\sim \text{Exp}(\lambda)$.

Minimize
\[
E|X-a|.
\]

Derivative condition gives
\[
P(X\le a)=\frac12.
\]

For exponential distribution:
\[
1-e^{-\lambda a}=\frac12
\Rightarrow a=\frac{\ln 2}{\lambda}.
\]

\textbf{Result:} $\boxed{a=\frac{\ln 2}{\lambda}}$.

\section*{Tutorial Question 4}

Battery lifetime is exponential with rate $\lambda_i$ with probability $p_i$.

We compute
\[
P(X>t+s\mid X>t).
\]

Using conditional probability:
\[
\frac{P(X>t+s)}{P(X>t)}.
\]

Mixture survival:
\[
P(X>t)=\sum_{i=1}^2 p_i e^{-\lambda_i t}.
\]

Thus
\[
P(X>t+s\mid X>t)
=\frac{\sum p_i e^{-\lambda_i(t+s)}}
{\sum p_i e^{-\lambda_i t}}.
\]

\section*{Tutorial Question 5}

$X\sim \text{Poisson}(\lambda_1)$,  
$Y\sim \text{Poisson}(\lambda_2)$ independent.

Given total $X+Y=n$, find $P(X=k\mid X+Y=n)$.

\[
P(X=k\mid X+Y=n)
=\frac{P(X=k,Y=n-k)}{P(X+Y=n)}.
\]

Substituting Poisson pmfs:
\[
=\binom{n}{k}
\left(\frac{\lambda_1}{\lambda_1+\lambda_2}\right)^k
\left(\frac{\lambda_2}{\lambda_1+\lambda_2}\right)^{n-k}.
\]

\textbf{Result:} Binomial$(n,\frac{\lambda_1}{\lambda_1+\lambda_2})$.

\section*{Tutorial Question 6}

$X\sim \text{Uniform}[-2,2]$, $Y=g(X)$ from graph.

\textbf{Procedure:}
\begin{itemize}
\item Identify intervals where $g(x)$ is monotone.
\item Express inverse mappings.
\item Use
\[
F_Y(y)=P(g(X)\le y).
\]
\item Differentiate to obtain
\[
f_Y(y)=\frac{d}{dy}F_Y(y).
\]
\end{itemize}

Graphs of $F_Y(y)$ and $f_Y(y)$ follow directly from piecewise mapping.

\section*{Tutorial Question 7}

Each period multiplies price by $u$ or $d$.

Let $N$ = number of upward moves in 1000 periods.

\[
N\sim \text{Binomial}(1000,0.52).
\]

Price after 1000 periods:
\[
S=u^N d^{1000-N}.
\]

Condition for $30\%$ increase:
\[
u^N d^{1000-N}\ge1.3.
\]

Take logarithms to obtain threshold value of $N$.

Approximate probability using normal approximation to binomial:
\[
N\approx N(1000p,1000p(1-p)).
\]

Compute
\[
P(N\ge N_0).
\]

\section*{Tutorial Question 8}

Repair time $X\sim \text{Exp}(1/2)$.

\subsection*{(a)}
\[
P(X>2)=e^{-(1/2)\cdot2}=e^{-1}.
\]

\subsection*{(b)}
\[
P(X\ge10\mid X\ge9)
=\frac{P(X\ge10)}{P(X\ge9)}
=\frac{e^{-5}}{e^{-4.5}}
=e^{-0.5}.
\]

\section*{Tutorial Question 9}

Let $Z=\sqrt{X^2+Y^2}$.

Joint density:
\[
f_{X,Y}(x,y)=f_X(x)f_Y(y).
\]

Use transformation to polar coordinates:
\[
x=r\cos\theta,\quad y=r\sin\theta.
\]

Jacobian $=r$.

Thus
\[
f_Z(z)=\int_0^{2\pi}
f_X(z\cos\theta)f_Y(z\sin\theta)\, z\, d\theta.
\]

\section*{Tutorial Question 10}

Let
\[
M_n=\max(X_1,\dots,X_n).
\]

Then
\[
P(M_n\le x)=F(x)^n.
\]

Joint distribution:
\[
P(M_n\le x, M_{n+1}\le y)
=
\begin{cases}
F(x)^nF(y), & x\le y,\\
F(y)^{n+1}, & y<x.
\end{cases}
\]

This follows because $M_{n+1}=\max(M_n,X_{n+1})$.

\end{document}